\documentclass[12pt]{article}

\usepackage{amsmath,amssymb}
\usepackage[a4paper,margin=1in]{geometry}

\title{CSE400 -- Fundamentals of Probability in Computing}
\author{Tutorial 1: Complete Lecture Scribe (Solved)}
\date{}

\begin{document}

\maketitle

\section*{Q1. Dividing 20 Dishes into Platters}

\textbf{Definitions}

\textbf{Combination:}
\[
{n \choose k} = \frac{n!}{k!(n-k)!}
\]

\textbf{Partitioning identical groups:}

If $n$ objects are divided into $k$ groups of size $r$ each and order does not matter:

\[
\text{Number of partitions} = \frac{n!}{(r!)^k k!}
\]

\subsection*{(a) Total number of ways to divide dishes}

Given:

\begin{itemize}
\item 20 distinct dishes
\item 4 groups
\item 5 dishes per group
\item groups are unlabeled
\end{itemize}

\[
N = \frac{20!}{(5!)^4 4!}
\]

\subsection*{(b) Probability each platter has one cuisine}

Assume:

\begin{itemize}
\item 5 Italian
\item 5 Indian
\item 5 Chinese
\item 5 Mexican
\end{itemize}

Total partitions:

\[
\frac{20!}{(5!)^4 4!}
\]

Only one arrangement satisfies the requirement.

\[
P = \frac{1}{\frac{20!}{(5!)^4 4!}}
\]

---
\newpage
\section*{Q2. Bus Waiting Time}

\textbf{Uniform Random Variable}

\[
X \sim U(a,b)
\]

\[
f(x) = \frac{1}{b-a}, \quad a \le x \le b
\]

Passenger arrival:

\[
X \sim U(0,30)
\]

Bus arrivals: 7:00, 7:15, 7:30.

\subsection*{(a) Wait less than 5 minutes}

Intervals:

\begin{itemize}
\item 7:10 -- 7:15
\item 7:25 -- 7:30
\end{itemize}

Total favorable time = 10 minutes.

\[
P = \frac{10}{30} = \frac{1}{3}
\]

\subsection*{(b) Wait more than 10 minutes}

Intervals:

\begin{itemize}
\item 7:00 -- 7:05
\item 7:15 -- 7:20
\end{itemize}

\[
P = \frac{10}{30} = \frac{1}{3}
\]

---
\newpage
\section*{Q3. Conditional Probability}

\[
P(A|B) = \frac{P(A \cap B)}{P(B)}
\]

Let:

\[
P(L) = 0.15, \quad P(E) = 0.25
\]

\[
P(L \cap E) = 0.08
\]

Union:

\[
P(L \cup E) = P(L) + P(E) - P(L \cap E)
\]

\[
= 0.15 + 0.25 - 0.08
\]

\[
= 0.32
\]

Complement:

\[
P(L^c \cap E^c) = 1 - 0.32 = 0.68
\]

\[
P(E^c) = 1 - 0.25 = 0.75
\]

\[
P(L^c | E^c) = \frac{0.68}{0.75} = 0.9067
\]

---
\newpage
\section*{Q4. Typographical Errors (Poisson)}

Poisson distribution:

\[
P(X = k) = \frac{e^{-\lambda} \lambda^k}{k!}
\]

Given:

\[
\lambda = 0.2
\]

\subsection*{(a)}

\[
P(X=0) = e^{-0.2}
\]

\[
P \approx 0.8187
\]

\subsection*{(b)}

\[
P(X \ge 2) = 1 - P(0) - P(1)
\]

\[
P(1) = 0.2 e^{-0.2}
\]

\[
P(X \ge 2) = 1 - e^{-0.2}(1 + 0.2)
\]

\[
P \approx 0.0176
\]

---
\newpage
\section*{Q5. Airplane Crashes}

Given:

\[
\lambda = 3.5
\]

\subsection*{(a) At least 2 accidents}

\[
P(X \ge 2) = 1 - P(0) - P(1)
\]

\[
P(0) = e^{-3.5}
\]

\[
P(1) = 3.5e^{-3.5}
\]

\[
P(X \ge 2) = 1 - e^{-3.5}(1+3.5)
\]

\subsection*{(b) At most 1 accident}

\[
P(X \le 1) = P(0) + P(1)
\]

\[
= e^{-3.5}(1+3.5)
\]

---
\newpage
\section*{Q6. Poisson Accidents on Highway}

\[
\lambda = 3
\]

\[
P(X \ge 3) = 1 - P(0) - P(1) - P(2)
\]

\[
P(0) = e^{-3}
\]

\[
P(1) = 3e^{-3}
\]

\[
P(2) = \frac{9}{2}e^{-3}
\]

\[
P(X \ge 3) = 1 - e^{-3}\left(1+3+\frac{9}{2}\right)
\]

\[
P \approx 0.5768
\]

Conditional probability:

\[
P(X \ge 3 | X \ge 1) =
\frac{P(X \ge 3)}{P(X \ge 1)}
\]

\[
P(X \ge 1) = 1 - P(0) = 1 - e^{-3}
\]

---
\newpage
\section*{Q7. Gaussian Random Variable}

Gaussian pdf:

\[
f_X(x) = \frac{1}{\sqrt{2\pi\sigma^2}} 
\exp\left(-\frac{(x-\mu)^2}{2\sigma^2}\right)
\]

Given:

\[
\mu = -3, \quad \sigma = 2
\]

Standardization:

\[
Z = \frac{X-\mu}{\sigma}
\]

Q-function:

\[
Q(x) = P(Z>x)
\]

\[
P(X \le 0) = 1 - Q(1.5)
\]

\[
P(X > 4) = Q(3.5)
\]

\[
P(|X+3|<2)
\]

\[
-5 < X < -1
\]

\[
P = \Phi(1) - \Phi(-1)
\]

\[
= 1 - 2Q(1)
\]

---
\newpage
\section*{Q8. Jury Decision}

Let:

\[
X \sim \text{Binomial}(12,\theta)
\]

Conviction requires at least 8 correct votes.

\[
P(\text{correct}) =
\sum_{k=8}^{12} {12 \choose k} \theta^k (1-\theta)^{12-k}
\]

---
\newpage
\section*{Q9. Given PMF}

\[
p(i) = c \frac{\lambda^i}{i!}, \quad i = 0,1,2,\dots
\]

Using normalization:

\[
\sum_{i=0}^{\infty} p(i) = 1
\]

\[
c \sum_{i=0}^{\infty} \frac{\lambda^i}{i!} = 1
\]

\[
\sum_{i=0}^{\infty} \frac{\lambda^i}{i!} = e^{\lambda}
\]

Thus:

\[
c = e^{-\lambda}
\]

---
\newpage
\section*{Q10. System Reliability}

\[
X \sim \text{Binomial}(n,p)
\]

System works if:

\[
X \ge \frac{n}{2}
\]

\subsection*{(a) Compare $n=3$ and $n=5$}

\[
P_3 = P(X \ge 2)
\]

\[
= 3p^2(1-p) + p^3
\]

\[
P_5 = P(X \ge 3)
\]

\[
= 10p^3(1-p)^2 + 5p^4(1-p) + p^5
\]

Solve:

\[
P_5 > P_3
\]

Result:

\[
p > \frac{1}{2}
\]

\subsection*{(b) General Case}

Compare systems of size $2k+1$ and $2k-1$.

The larger system performs better when

\[
p > \frac{1}{2}
\]

\end{document}